\chapter{Typhoid (Enteric Fever) Test}

\section*{What Is This Test?}
Typhoid testing evaluates suspected enteric fever caused by \textit{Salmonella enterica} serovar Typhi; clinically similar paratyphoid fever is caused by \textit{S. Paratyphi}. Laboratory confirmation is important because typhoid can resemble malaria, dengue, viral infection, and other causes of prolonged fever. The preferred diagnostic test is culture of blood collected before antibiotics when possible. Stool culture, urine culture, bone-marrow culture, and molecular tests may be used in selected situations. The Widal test and other antibody-based rapid tests are not reliable enough to confirm acute typhoid on their own.

\section*{Alternative Names}
Typhoid Test; Enteric Fever Test; \textit{Salmonella Typhi} Culture; Typhoid Blood Culture; Typhi O and H Antibody Test; Widal Test; Typhoid PCR

\section*{Principle}
Blood culture uses an automated broth system to detect viable \textit{Salmonella} organisms, followed by identification and antimicrobial susceptibility testing. Multiple blood cultures may be needed because sensitivity is limited and previous antibiotics reduce recovery. Bone-marrow culture is more invasive but can remain positive after antibiotic exposure and has higher sensitivity in some studies. Stool and urine cultures are less sensitive for initial diagnosis, particularly early in illness. Widal testing measures antibodies to O and H antigens; antibody titers are affected by previous infection, vaccination, background exposure, and cross-reacting infections, so a single result cannot reliably distinguish acute from past infection.

\section*{History}
The Widal agglutination test was developed in the late nineteenth century and became widely used because it was inexpensive and technically simple. Modern clinical microbiology has established culture as the preferred method for confirming enteric fever and for obtaining an isolate for antimicrobial susceptibility testing. Molecular and rapid serologic tests continue to be evaluated, but their performance varies by assay and setting.

\section*{Why This Test Is Ordered}
\begin{itemize}
  \item Persistent or progressively increasing fever with headache, malaise, abdominal symptoms, diarrhea or constipation, or relative bradycardia when clinically relevant
  \item Fever after travel or residence in an area where typhoid or paratyphoid is endemic, or after a recognized exposure
  \item Suspected enteric fever in a patient whose symptoms are not explained by malaria, dengue, urinary infection, pneumonia, or another febrile illness
  \item Confirmation of infection and antimicrobial susceptibility testing before selecting or narrowing antibiotic therapy
  \item Investigation of a possible outbreak or public-health exposure
\end{itemize}

\section*{Primary vs Secondary Test}
Blood culture is the primary test for laboratory confirmation of acute typhoid or paratyphoid fever. Two or more cultures may be collected from separate venipunctures when suspicion is high. Stool culture is an adjunct and may be useful later in illness or for detecting fecal shedding. Bone-marrow culture is a specialist test when blood cultures are negative but suspicion remains high. Widal and other antibody tests are secondary, low-specificity tests and should not be used as the sole basis for diagnosis or treatment.

\section*{How This Test Is Performed}
For blood culture, the skin is disinfected carefully and blood is collected into one or more aerobic culture bottles, ideally before the first antibiotic dose. The required volume depends on age and the laboratory system; collection should follow the laboratory's instructions. Stool is collected in a clean, leak-proof container without urine or disinfectant contamination. A laboratory may perform antimicrobial susceptibility testing on any isolate recovered. PCR or other rapid assays require the specimen type and transport conditions specified by the performing laboratory.

\section*{What Preparation Is Needed}
No fasting is required. Tell the clinician about recent or current antibiotics, travel, vaccination, prior typhoid infection, and immunosuppression. Antibiotics taken before collection can cause a false-negative culture. Do not delay urgent clinical care solely to obtain a test result.

\section*{Contraindications and Precautions}
\begin{itemize}
  \item There are no specific contraindications to blood or stool culture beyond the usual precautions for venipuncture and specimen collection. A negative blood culture does not exclude typhoid, especially after antibiotics, with low bacterial burden, or when too little blood was collected. A positive Widal result does not by itself establish acute infection. Suspected cases should be managed according to local public-health and antimicrobial-resistance guidance; confirmed \textit{Salmonella Typhi} or \textit{S. Paratyphi} infection may be notifiable.
\end{itemize}
\section*{Risks}
Blood collection may cause brief pain, bruising, bleeding, or rarely infection. Stool collection has no significant physical risk. The principal clinical risk is misdiagnosis from an unreliable serologic result or failure to recognize another serious cause of fever.

\section*{What Happens After the Test}
A preliminary positive blood-culture signal is communicated urgently for organism identification and susceptibility testing. Final culture reporting may require several days; a negative result may take longer according to the laboratory system. Patients with confirmed or strongly suspected enteric fever require clinician-directed treatment and follow-up. Persistent fever, severe abdominal pain, gastrointestinal bleeding, confusion, dehydration, or worsening illness requires urgent medical assessment.

\section*{Understanding Results}

\subsection*{Negative}
\begin{itemize}
  \item \textbf{Blood culture negative:} No \textit{Salmonella} was recovered from that specimen. This does not exclude typhoid, particularly after antibiotics or when only one culture was obtained.
  \item \textbf{Stool culture negative:} No organism was recovered from the stool specimen; early or intermittent shedding can produce a negative result.
  \item \textbf{Widal negative:} Antibody levels were below the laboratory cutoff. Early infection may still be possible, and the result should not override clinical assessment and culture testing.
\end{itemize}

\subsection*{Positive}
\begin{itemize}
  \item \textbf{Blood culture positive for \textit{S. Typhi} or \textit{S. Paratyphi}:} Confirms enteric fever, subject to laboratory identification and contamination assessment. Antimicrobial susceptibility testing should guide treatment.
  \item \textbf{Stool culture positive:} Supports infection or intestinal shedding, but interpretation depends on symptoms, specimen timing, and whether the isolate is confirmed as a pathogenic serovar.
  \item \textbf{Widal positive:} Indicates antibodies at or above the local cutoff, not proof of acute typhoid. Results must be interpreted with illness duration, local baseline titers, vaccination, exposure history, and culture or molecular testing.
\end{itemize}

\section*{Normal Results}
Blood culture: \textbf{No \textit{Salmonella} isolated}. Stool culture: \textbf{No \textit{S. Typhi} or \textit{S. Paratyphi} isolated}. Widal: \textbf{Not detected or below the laboratory cutoff}; this does not independently rule out early disease.

\section*{Follow-up Tests}
\begin{itemize}
  \item \textbf{Antimicrobial susceptibility testing:} Required for a recovered isolate because resistance patterns vary geographically and over time.
  \item \textbf{Repeat blood cultures:} Considered when initial cultures are negative but suspicion remains high, especially before antibiotics whenever feasible.
  \item \textbf{Malaria testing and dengue or other febrile-illness evaluation:} Selected according to travel history, season, geography, symptoms, and local epidemiology.
  \item \textbf{Complete blood count, liver tests, electrolytes, and renal tests:} Used to assess illness severity and complications; these tests do not confirm typhoid.
\end{itemize}

\section*{References}
\begin{enumerate}
  \item Centers for Disease Control and Prevention. Clinical Guidance for Typhoid Fever and Paratyphoid Fever. Updated 2024.
  \item Centers for Disease Control and Prevention. Yellow Book: Typhoid and Paratyphoid Fever. Updated 2024.
  \item World Health Organization. WHO laboratory manual for the diagnosis of typhoid and paratyphoid fever. WHO.
\end{enumerate}

\clearpage
\section*{Regulatory Status and Authorization Date}
This chapter describes a test category, not a specific commercial product. FDA, CDSCO, EU IVDR, and MHRA approval or registration dates therefore cannot be assigned without the assay manufacturer, product name, instrument, and intended use. For laboratory-developed tests, an individual agency approval date may not exist. See the Preface for the interpretation of regulatory status.

\clearpage
